\section{Implementasi dalam C}

Berikut implementasi Caesar Cipher dan Fast Modular Exponentiation dalam C.

\subsection{Caesar Cipher}
\begin{lstlisting}[style=cstyle,caption={Implementasi Caesar Cipher dalam C}]
#include <stdio.h>
#include <ctype.h>
#include <string.h>

void caesar_encrypt(char *plaintext, int shift) {
    int i;
    for (i = 0; plaintext[i] != '\0'; i++) {
        if (isalpha(plaintext[i])) {
            char base = isupper(plaintext[i]) ? 'A' : 'a';
            plaintext[i] = base + (plaintext[i] - base + shift) % 26;
        }
    }
}

void caesar_decrypt(char *ciphertext, int shift) {
    caesar_encrypt(ciphertext, 26 - shift);
}

int main() {
    char msg[] = "HELLO CRYPTO";
    caesar_encrypt(msg, 3);
    printf("Encrypted: %s\n", msg);
    caesar_decrypt(msg, 3);
    printf("Decrypted: %s\n", msg);
    return 0;
}
\end{lstlisting}

\subsection{Fast Modular Exponentiation}
\begin{lstlisting}[style=cstyle,caption={Fast Modular Exponentiation dalam C}]
#include <stdio.h>

long long mod_pow(long long base, long long exp, long long mod) {
    long long result = 1;
    base = base % mod;
    while (exp > 0) {
        if (exp % 2 == 1)
            result = (result * base) % mod;
        exp = exp / 2;
        base = (base * base) % mod;
    }
    return result;
}

int main() {
    printf("2^10 mod 1000 = %lld\n", mod_pow(2, 10, 1000));
    printf("3^644 mod 645 (RSA) = %lld\n", mod_pow(3, 644, 645));
    return 0;
}
\end{lstlisting}
